\documentclass[11pt,a4paper]{article}
\usepackage[left=23mm,right=23mm,top=22mm,bottom=22mm,footskip=11mm]{geometry}
\usepackage{amsmath,amssymb}
\usepackage[no-math]{fontspec}
\usepackage{xeCJK}
\setmainfont{Latin Modern Roman}
\setCJKmainfont{FandolSong-Regular.otf}[Path=assets/fonts/fandol/,BoldFont=FandolSong-Bold.otf,ItalicFont=FandolSong-Regular.otf]
\setCJKsansfont{FandolSong-Regular.otf}[Path=assets/fonts/fandol/,BoldFont=FandolSong-Bold.otf]
\renewcommand{\bfdefault}{b}
\xeCJKsetup{PunctStyle=quanjiao,CheckSingle=true}
\usepackage{setspace}\setstretch{1.10}
\setlength{\parindent}{2em}\setlength{\parskip}{4pt}
\usepackage{xcolor,titlesec,tcolorbox,tikz,enumitem,longtable,booktabs,array,calc,needspace}
\tcbuselibrary{breakable,skins}\usetikzlibrary{arrows.meta,positioning}
\definecolor{sitepaper}{HTML}{FFFEFC}\pagecolor{sitepaper}
\definecolor{navy}{HTML}{10136B}\definecolor{inkgray}{HTML}{474747}\definecolor{warm}{HTML}{94662C}\definecolor{teal}{HTML}{236B66}
\titleformat{\section}{\Large\bfseries\color{navy}}{}{0pt}{}
\titleformat{\subsection}{\normalsize\bfseries\color{navy}}{}{0pt}{}
\titlespacing*{\section}{0pt}{17pt}{9pt}\titlespacing*{\subsection}{0pt}{12pt}{6pt}
\newtcolorbox{sboxenv}[2]{enhanced,breakable,colframe=#1,colback=#1!3!white,colbacktitle=#1,coltitle=white,fonttitle=\bfseries,title={#2},boxrule=0.45pt,arc=1.4pt,left=8pt,right=8pt,top=6pt,bottom=6pt,before skip=9pt,after skip=9pt,pad at break*=2mm}
\setlist[enumerate]{leftmargin=2.2em,itemsep=3pt,parsep=0pt,topsep=3pt}
\setlist[itemize]{label={},leftmargin=1em,itemsep=3pt}
\providecommand{\tightlist}{\setlength{\itemsep}{3pt}\setlength{\parskip}{0pt}}
\usepackage[unicode,hidelinks,linktoc=all]{hyperref}\usepackage{bookmark}
\hypersetup{pdftitle={丢番图逼近的初等方法}}
\renewcommand{\contentsname}{目录}\setcounter{tocdepth}{1}\setcounter{secnumdepth}{0}
\widowpenalty=10000\clubpenalty=10000\displaywidowpenalty=10000\predisplaypenalty=5000
\allowdisplaybreaks[1]\raggedbottom\emergencystretch=2em


\usepackage{fancyhdr}
\pagestyle{fancy}\fancyhf{}\setlength{\headheight}{14pt}
\fancyhead[L]{\small\color{inkgray} 数学手记}
\fancyhead[R]{\small\color{inkgray} 丢番图逼近的初等方法}
\fancyfoot[C]{\small\thepage}
\renewcommand{\headrulewidth}{0pt}
\newcommand{\articleheading}[1]{\section{#1}}
\begin{document}
\thispagestyle{plain}
\begin{center}{\LARGE\bfseries 丢番图逼近的初等方法\par}\vspace{5pt}
{\large 从逼近的方向到分母的平方阶\par}\end{center}

已知 $5\sqrt 2$ 比整数 $7$ 大一点。能否只凭这个信息，找一个正整数倍的 $\sqrt 2$，使它比某个整数小一点？若还要求分母为 $q$ 的逼近分数与 $\sqrt 2$ 相差不到 $\frac{1}{q^2}$，同一个办法还管用吗？

这两个问题的区别在于，放大倍数可以换取所需的方向，却也放大了分母。我们先在单位圆上寻找小误差，再用一对有理端点同时记住误差的方向和分母。这些证明都不使用连分数。

\articleheading{§1 单位圆上最近的两个点}

记 $\lfloor x\rfloor$ 为不超过 $x$ 的最大整数，$\{x\}=x-\lfloor x\rfloor$ 为小数部分，$\|x\|=\min_{m\in\mathbb Z}|x-m|$ 为到最近整数的距离。对 $q>0$，$q\alpha-p>0$ 表示 $q\alpha$ 在整数 $p$ 的右侧，同时分数 $\frac pq$ 从下方逼近 $\alpha$。这两种说法描述的是同一个符号。

\begin{sboxenv}{navy}{定理 1　Dirichlet 有限逼近}

给定实数 $\alpha$ 和正整数 $N$，存在整数 $p,q$，满足

\[
1\le q\le N,\qquad |q\alpha-p|\le\frac{1}{N+1}.\qquad\mathrm{(1)}
\]

当 $\alpha$ 无理时，右端可以取严格不等号。
\end{sboxenv}

\textbf{想法}　把 $0,\{\alpha\},\ldots,\{N\alpha\}$ 放入 $N$ 个等长的半开区间，两个落在同一格的点相减，便得到小于 $\frac1N$ 的误差。要得到（1）中的 $N+1$，可以把 $0$ 与 $1$ 相接，把首尾两点之间的空隙也算进去。

\textbf{证明}　若有两点重合，取它们下标之差的绝对值为 $q$，便有 $q\alpha\in\mathbb Z$。否则，按圆周顺序排列这 $N+1$ 个点。相邻弧长之和为 $1$，最短弧长 $h$ 至多为 $\frac{1}{N+1}$。设该弧两端下标为 $i,j$。沿弧走过的距离与 $(j-i)\alpha$ 相差一个整数，跨过原点时也如此。令 $q=|j-i|$，便有某个整数 $p$ 使 $|q\alpha-p|=h$。若 $\alpha$ 无理，等号 $h=\frac1{N+1}$ 会使 $q\alpha=p\pm\frac1{N+1}$ 为有理数，矛盾。

这个界对所有实数统一成立时已经最优。取 $\alpha=\frac1{N+1}$，圆上的点恰好等分圆周，而 $1\le q\le N$ 时 $\|q\alpha\|$ 的最小值正是 $\frac1{N+1}$。

\begin{sboxenv}{navy}{推论 1　平方阶有理逼近}

每个无理数 $\alpha$ 都有无穷多个既约分数 $\frac pq$ 满足

\[
0<\left|\alpha-\frac pq\right|<\frac1{q^2}.\qquad\mathrm{(2)}
\]
\end{sboxenv}

\textbf{证明}　将（1）除以 $q$，所得上界为 $\frac1{q(N+1)}<\frac1{q^2}$。若原分数约成 $\frac ab$，则 $b\le q$，故它也满足以 $b$ 为分母的（2）。使 $N$ 趋于无穷，构造出的分数与 $\alpha$ 的距离至多为 $\frac1{N+1}$，从而趋于零。有限多个有理数到无理数 $\alpha$ 的距离有正的最小值，故这些既约分数必有无穷多个。

\clearpage\articleheading{§2 用正整数倍改变误差方向}

\begin{sboxenv}{navy}{引理 2　正倍数换边}

设 $q\alpha-p=\sigma\delta$，其中 $q$ 是正整数，$p$ 是整数，$\sigma\in\{-1,1\}$，$0<\delta<1$，且 $\alpha$ 无理。取 $m=\lfloor\frac1\delta\rfloor$，则

\[
(mq)\alpha-(mp+\sigma)=-\sigma(1-m\delta),
\qquad 0<1-m\delta<\delta.\qquad\mathrm{(3)}
\]
\end{sboxenv}

\textbf{想法}　每走 $q$ 步，小数部分就沿同一方向移动 $\delta$。走到快满一圈时，从下一个整数回看，误差自然换了符号。这个办法把下标始终保留为正数。

\textbf{证明}　$\delta=|q\alpha-p|$ 无理，故 $\frac1\delta$ 也无理。从 $m<\frac1\delta<m+1$ 得 $m\delta<1<(m+1)\delta$，即 $0<1-m\delta<\delta$。把已知等式乘以 $m$，再从两边减去 $\sigma$，即得（3）。

结合定理 1，对任意 $\varepsilon>0$，先选取 $0<|q\alpha-p|<\min(\varepsilon,\frac12)$。保留这个解，再用（3）构造另一个解，便在整数左右各有一个误差小于 $\varepsilon$ 的正下标。令 $\varepsilon$ 趋零，每一侧的下标均无界，因为任意有限多个正下标对应的 $\|q\alpha\|$ 都有正的最小值。

\begin{sboxenv}{warm}{例 1　五分之七与七十分之九十九}

取 $\alpha=\sqrt2$。由 $7^2<50$ 知 $\delta=5\sqrt2-7>0$。有理化给出 $\frac1\delta=5\sqrt2+7$，它介于 $14$ 与 $15$ 之间，故 $m=14$。
\end{sboxenv}

代入（3）得到 $70\sqrt2-99=-(1-14\delta)<0$。还可以用平方差精确核算两次误差：

\[
5\sqrt2-7=\frac1{5\sqrt2+7},\qquad
99-70\sqrt2=\frac1{99+70\sqrt2}.
\]

于是 $\frac75<\sqrt2<\frac{99}{70}$。这个例子中新分数尤其精确，但一般的（3）只估计 $1-m\delta<\delta$，没有给出 $1-m\delta<\frac1{mq}$。要对新分母也取得平方阶，须在构造中同时跟踪分母。

\clearpage\articleheading{§3 从上下两侧各取无穷多个平方阶逼近}

\begin{sboxenv}{navy}{定理 3　指定方向的平方阶逼近}

对每个无理数 $\alpha$，以下两个不等式各有无穷多个既约分数解，且分母可以任意大：

\[
0<\alpha-\frac ab<\frac1{b^2},\qquad
0<\frac cd-\alpha<\frac1{d^2}.\qquad\mathrm{(4)}
\]
\end{sboxenv}

\textbf{想法}　若用 $\frac ab<\alpha<\frac cd$ 夹住 $\alpha$，两端的方向就固定了。特别希望 $bc-ad=1$，因为此时区间宽度恰为 $\frac1{bd}$。只要右端分母 $d$ 超过左端分母 $b$，左端就自动符合（4）。问题归结为怎样缩小区间，同时保住行列式 $1$。

\textbf{证明}　令 $k=\lfloor\alpha\rfloor$，从端点 $\frac k1,\frac{k+1}1$ 出发。对当前端点作中介分数 $u=\frac{a+c}{b+d}$。直接通分得

\[
u-\frac ab=\frac{1}{b(b+d)},\qquad
\frac cd-u=\frac{1}{d(b+d)}.\qquad\mathrm{(5)}
\]

因 $\alpha$ 无理，必有 $u\ne\alpha$。留下含 $\alpha$ 的半边，即在 $u<\alpha$ 时更新左端，在 $u>\alpha$ 时更新右端。新端点对的行列式仍为 $1$，因为 $b(a+c)-a(b+d)=bc-ad$，而 $(b+d)c-(a+c)d=bc-ad$。新分数的分子与分母的任何公因数都整除其中一个行列式，故它已既约。由此（5）可以无限次使用。

两端都会更新无穷多次。若从某时起左端固定为 $\frac ab$，右端分母每次增加 $b$，故趋于无穷，而夹逼区间的宽 $\frac1{bd}$ 趋于零。这迫使 $\alpha=\frac ab$，矛盾。若右端最终固定，交换左右可得同样的矛盾。

现在看一段连续的左端更新结束、刚刚更新右端的时刻。右端的新分母为两旧分母之和，大于当前左分母 $b$。因而

\[
0<\alpha-\frac ab<\frac1{bd}<\frac1{b^2}.
\]

两端都无限次更新，故这种由左换右的时刻有无穷多个，相邻两次之间左分母严格增加。于是下方有无穷多个不同的解。考察由右换左的时刻，则 $b>d$，从而 $0<\frac cd-\alpha<\frac1{bd}<\frac1{d^2}$，上方亦然。

\textbf{把构造算一遍}　对 $\sqrt2$，从 $1,2$ 出发，依次插入 $\frac32,\frac43,\frac75,\frac{10}{7}$。它们相对于 $\sqrt2$ 依次在上、下、下、上方，可逐个比较分子的平方与分母平方的两倍。最后留下的区间是 $(\frac75,\frac{10}{7})$。即使不知道 $\sqrt2$ 的任何小数位，也已经知道 $0<\sqrt2-\frac75<\frac1{35}<\frac1{25}$。这个保证恰在更新另一端时获得。

\begin{sboxenv}{inkgray}{辨析　给定分母上限时，方向会改变结论}

固定 $N\ge2$，取无理数 $0<\alpha<\frac1{N(N+1)}$。对每个 $1\le q\le N$，有 $0<q\alpha<\frac1{N+1}$；而任何位于 $q\alpha$ 上方的整数与它的距离都至少为 $1-q\alpha>1-\frac1{N+1}$。因此（1）不能在固定 $N$ 时任意指定方向。（4）允许分母随方向而选。
\end{sboxenv}

\clearpage\articleheading{§4 从逼近整数到命中任意区间}

\begin{sboxenv}{navy}{定理 4　无理旋转的每条尾轨道都稠密}

若 $\alpha$ 无理，则对任意实数 $\beta$、非负整数 $L$ 及 $0\le u<v\le1$，存在整数 $n>L$ 使 $u<\{\beta+n\alpha\}<v$。因此每个这样的区间都被命中无穷多次。
\end{sboxenv}

\textbf{想法}　§2 已经给出任意小的正步长。把这个步长沿圆周重复使用，到快满一圈时，整圆上留下的每段空隙都很小。平移这一组点，就能从任意起点出发。

\textbf{证明}　选 $q>0$ 使 $\delta=\{q\alpha\}$ 满足 $0<\delta<v-u$。取 $m=\lfloor\frac1\delta\rfloor$，则 $0,\delta,\ldots,m\delta$ 是圆上一组严格递增的点，相邻间隙长 $\delta$，最后跨过原点的间隙长 $1-m\delta<\delta$。把整组点同时平移 $\beta+(L+1)\alpha$ 并取小数部分，圆周上的最大间隙仍至多为 $\delta$。开区间 $(u,v)$ 的长度大于 $\delta$，不可能完全容于相邻两点间的空隙，故至少包含一个平移后的点。它等于 $\{\beta+(L+1+jq)\alpha\}$，其中 $0\le j\le m$。取 $n=L+1+jq$ 即得。由于 $L$ 任意，命中次数必为无穷。

\begin{sboxenv}{warm}{例 2　哪些二的幂以前四位 2026 开头？}

证明有无穷多个正整数 $n$，使 $2^n$ 的十进制表示以前四位 $2026$ 开头。
\end{sboxenv}

\textbf{解}　设 $\alpha=\log_{10}2$。若 $\alpha=\frac rs$，取正整数 $r,s$，则 $10^r=2^s$，左边含有质因子 $5$，右边没有，矛盾。写 $n\alpha=k+t$，其中 $k=\lfloor n\alpha\rfloor$，$t=\{n\alpha\}$，于是 $2^n=10^k\,10^t$。当 $k\ge3$ 且 $2.026<10^t<2.027$ 时，

\[
2026\cdot10^{k-3}<2^n<2027\cdot10^{k-3},
\]

所以前四位必为 $2026$。由定理 4，无穷多个任意大的 $n$ 使 $t$ 落入 $(\log_{10}2.026,\log_{10}2.027)$。同一证明也适用于任何预先指定、首位非零的有限数字串。

\textbf{若下标也有限制}　固定正整数 $h$ 与余数 $0\le r<h$。因 $h\alpha$ 无理，在定理 4 中将步长换成 $h\alpha$，起点换成 $r\alpha$，便知 $\{(r+jh)\alpha\}$ 仍无穷多次命中每个开区间。例如在例 2 中，若额外要求指数 $n$ 除以 $6$ 余 $5$，同样得到无穷多个以前四位 $2026$ 开头的二的幂。

\clearpage\articleheading{§5 用整数下界判断逼近能有多好}

\begin{sboxenv}{navy}{命题 5　根号二的平方阶不能提高}

对任意整数 $p$ 和正整数 $q$，若 $|\sqrt2-\frac pq|<1$，则

\[
\left|\sqrt2-\frac pq\right|>\frac{1}{(2\sqrt2+1)q^2}.\qquad\mathrm{(6)}
\]
\end{sboxenv}

\textbf{想法}　把小量乘以共轭因子。原来的差可以很小，分子里出现的非零整数却至少有绝对值 $1$。

\textbf{证明}　$2q^2-p^2$ 是非零整数，否则 $\sqrt2$ 为有理数（既约的 $p^2=2q^2$ 会推出 $p,q$ 都是偶数）。由假设 $0<\frac pq<\sqrt2+1$，故

\[
\begin{aligned}
\left|\sqrt2-\frac pq\right|
&=\frac{|2q^2-p^2|}{q^2(\sqrt2+\frac pq)}\\
&\ge\frac1{q^2(\sqrt2+\frac pq)}
{}>\frac{1}{(2\sqrt2+1)q^2}.
\end{aligned}
\]

与（4）合看，对 $\sqrt2$ 而言，分母的二次方给出了正确的阶。特别地，对任意固定 $\eta>0$，当 $q$ 足够大时，误差小于 $\frac1{q^{2+\eta}}$ 已不可能。

\textbf{怎样找到好分数？}　要使上式的分子等于 $1$，可把 §2 例子中的平方差推广成恒等式。对 $n\ge1$ 写

\[
(1+\sqrt2)^n=p_n+q_n\sqrt2,\qquad
(1-\sqrt2)^n=p_n-q_n\sqrt2.
\]

由二项式定理，$p_n,q_n$ 为正整数。相乘得 $p_n^2-2q_n^2=(-1)^n$，故 $p_n,q_n$ 互素。相减又得 $2\sqrt2\,q_n=(1+\sqrt2)^n-(1-\sqrt2)^n$，所以 $q_n\to\infty$。由于 $|1-\sqrt2|<1$，第二个恒等式给出 $\frac{p_n}{q_n}\to\sqrt2$。分子的符号随 $n$ 交替，且

\[
q_n^2\left|\sqrt2-\frac{p_n}{q_n}\right|
=\frac1{\sqrt2+\frac{p_n}{q_n}}
\longrightarrow\frac1{2\sqrt2}.
\]

$\frac11,\frac32,\frac75,\frac{17}{12},\frac{41}{29},\frac{99}{70}$ 正是前六项。对任何别的趋于 $\sqrt2$ 的有理数列，平方差的绝对值仍至少为 $1$，所以归一化误差 $q^2|\sqrt2-\frac pq|$ 的下极限也至少为 $\frac1{2\sqrt2}$。这说明平方阶的最佳渐近常数是 $\frac1{2\sqrt2}$。

\clearpage\articleheading{§6 从多项式的分母到一个超越数}

\begin{sboxenv}{navy}{定理 6　Liouville 的初等下界}

设实数 $\alpha$ 是次数 $d\ge2$ 的代数数，即它在有理数域上的最小多项式次数为 $d$。存在只依赖 $\alpha$ 的常数 $C>0$，使任意整数 $p$、正整数 $q$ 都满足

\[
\left|\alpha-\frac pq\right|\ge\frac C{q^d}.\qquad\mathrm{(7)}
\]
\end{sboxenv}

\textbf{想法}　§5 用的是 $x^2-2$。把它换成 $\alpha$ 的最小多项式，代入 $\frac pq$ 的分母最多是 $q^d$。剩下的问题是怎样把多项式值的下界转化为自变量距离的下界。

\textbf{证明}　将最小多项式清去系数的分母，得整系数多项式 $P$，仍在有理数域上不可约。因 $d\ge2$，它没有有理根，故 $q^dP(\frac pq)$ 是非零整数，从而 $|P(\frac pq)|\ge q^{-d}$。令 $M=\max_{|x-\alpha|\le1}|P'(x)|$。导数连续且不恒为零，故 $0<M<\infty$。若 $|\alpha-\frac pq|<1$，中值定理给出

\[
q^{-d}\le |P(\tfrac pq)-P(\alpha)|\le M\left|\tfrac pq-\alpha\right|.
\]

若距离至少为 $1$，它自然也至少为 $q^{-d}$。取 $C=\min(1,\frac1M)$，即得（7）。

\begin{sboxenv}{warm}{例 3　把逼近速度写进小数位}

数 $\xi=\sum_{j=1}^{\infty}10^{-j!}$ 是超越数。
\end{sboxenv}

\textbf{解}　该级数由几何级数控制，因而收敛。截断到第 $m$ 项，写成 $\frac{P_m}{Q_m}$，其中 $Q_m=10^{m!}$。误差从第 $(m+1)!$ 位才开始。因后续指数至少逐次增加 $1$，

\[
0<\xi-\frac{P_m}{Q_m}
\le\sum_{k=(m+1)!}^{\infty}10^{-k}
=\frac{10}{9}Q_m^{-(m+1)}
<2Q_m^{-(m+1)}.\qquad\mathrm{(8)}
\]

先排除有理数。若 $\xi=\frac ab$，$b>0$，则严格为正的差至少为 $\frac1{bQ_m}$，与（8）在 $m$ 很大时矛盾。再假设 $\xi$ 是次数 $d\ge2$ 的代数数。（7）和（8）将推出 $C<2Q_m^{d-m-1}$。当 $m\to\infty$ 时右边趋于零，矛盾。于是 $\xi$ 既非有理数，也非代数无理数，只能是超越数。

\clearpage\articleheading{§7 同一个分母逼近多个实数}

\begin{sboxenv}{navy}{定理 7　Dirichlet 同步逼近}

给定实数 $\alpha_1,\ldots,\alpha_s$ 和正整数 $H$，存在同一个整数 $1\le q\le H^s$ 以及整数 $p_1,\ldots,p_s$，使

\[
|q\alpha_i-p_i|<\frac1H\qquad(1\le i\le s).\qquad\mathrm{(9)}
\]

若至少一个坐标无理，则有无穷多个不同的 $q$ 使所有坐标同时满足 $|\alpha_i-\frac{p_i}{q}|<q^{-1-\frac1s}$。
\end{sboxenv}

\textbf{想法}　对每个坐标分别使用定理 1 会得到不同的分母。把各个小数部分放进同一个向量，两点所在的小盒同时控制了每个坐标，而两个下标之差只有一个。

\textbf{证明}　把单位半开立方体 $[0,1)^s$ 的每条边等分为 $H$ 段，就有 $H^s$ 个半开小立方体。放入 $H^s+1$ 个点

\[
x_j=(\{j\alpha_1\},\ldots,\{j\alpha_s\}),\qquad 0\le j\le H^s .
\]

由鸽巢原理，有 $x_j,x_k$ 同处一格，设 $j>k$。取 $q=j-k$ 及 $p_i=\lfloor j\alpha_i\rfloor-\lfloor k\alpha_i\rfloor$。则 $q\alpha_i-p_i$ 正是同格两点第 $i$ 坐标之差，绝对值严格小于 $\frac1H$。又 $q\le H^s$，故 $H\ge q^{\frac1s}$。将（9）除以 $q$ 即得所述指数。若某一坐标 $\alpha_i$ 无理，则每个固定 $q$ 的 $\|q\alpha_i\|$ 都大于零，有限多个 $q$ 的这些距离有正的最小值。因此 $H$ 任意增大时，分母必取无穷多个不同值。

\begin{sboxenv}{warm}{例 4　用分母七同时逼近两个平方根}

对 $(\sqrt2,\sqrt3)$ 取 $H=3$。$x_2$ 与 $x_9$ 的第一坐标都在 $[\frac23,1)$，第二坐标都在 $[\frac13,\frac23)$。相减所得的公共分母为 $9-2=7$。
\end{sboxenv}

\textbf{解}　第一坐标的整数部分之差为 $12-2=10$，第二坐标为 $15-3=12$。算出的两个误差可以直接核实：

\[
0<10-7\sqrt2=\frac2{10+7\sqrt2}<\frac13,\qquad
0<7\sqrt3-12=\frac3{7\sqrt3+12}<\frac13.
\]

所以 $\frac{10}{7}$ 从上方逼近 $\sqrt2$，$\frac{12}{7}$ 从下方逼近 $\sqrt3$。两个方向恰好相反，取同一格所保证的是每个误差的绝对值。

回到开头的问题。若题目指定整数的哪一侧，先看是否限制了分母：只求误差趋零，可以用（3）换边；还要求平方阶，就用（5）保存两端分母的乘积。若想证明某种精度不可能，则反过来找一个非零整数，如 $2q^2-p^2$ 或 $q^dP(\frac pq)$。上界来自构造，下界来自整数之间的间隔。
\end{document}
